\begin{proof}
Take $\alpha=e_1$, $\beta=(1,0,1)^\top$. The only nonzero rows and columns are indexed by $(\Seps,a,b,aa)$, and the entire nonzero operator is
\[
 H=\begin{pmatrix}1&0&1&1\\0&1&0&0\\1&0&0&0\\1&0&0&0\end{pmatrix}.
\]
Its eigenvectors $(2,0,1,1)/\sqrt6$, $(0,1,0,0)$ and $(1,0,-1,-1)/\sqrt3$ have eigenvalues $2,1,-1$. Thus the rank is three and the realization is minimal.

For a complex rank-one candidate with parameters $(p,q)$, put $x=|p|^2$, $t=|q|$, $m=1-x-t^2>0$, and $z=q+p^2$. Projecting $H$ away from its possible output range gives the finite matrix
$M=H^2-mww^*$, $w=(1+z,p,1,1)^\top$. On $\Sran H$ one eigenvalue is one; the other two, $0\le\mu_-\le1\le\mu_+\le4$, satisfy
\[
 Q_0(\lambda)=(\lambda-4)(\lambda-1)+4A(\lambda-1)+B(\lambda-4),
\]
where $A=m|2+z|^2/6$ and $B=m(x+|1-z|^2/3)$. Direct evaluation gives
\[
 Q_0(4/3)=\tfrac23\{m(4\Re z-|z|^2-4x)-4/3\}.
\]
Since $t+x<1+t-t^2\le5/4$, the reverse triangle inequality yields
$4\Re z-|z|^2-4x\le4t-(t+x)^2$. If the right side is negative the next bound is immediate; otherwise both factors decrease with $x$. Therefore
$m(4\Re z-|z|^2-4x)\le(1-t^2)(4t-t^2)$.
For $y=t-1/2\in[-1/2,1/2]$, the exact identity
\[
 4-3(1-t^2)(4t-t^2)
 =\tfrac{51}{4}(y-1/17)^2+\tfrac5{272}
   +3y^2(y+1/2)(5/2-y)
\]
makes $Q_0(4/3)\le-5/1224$. Hence
$\mu_+\ge4/3+(5/1224)/(4/3)=727/544$, proving the lower bound. Subsection~\ref{I:S:app:commupper} gives a rational candidate and its strict upper certificate.

For the comparison put $r=2^{1/3}$ and let $h(\Seps)=2-r$, $h(a)=r^2-1$, $h(aa)=r$, with all other coefficients zero. Its unary block
\[
 K=\begin{pmatrix}2-r&r^2-1&r\\r^2-1&r&0\\r&0&0\end{pmatrix}
\]
has characteristic polynomial $(\lambda-2)(\lambda-1)(\lambda+1)$. Let $b_0=r^{-1}$ and $h_*(a^n)=(2-r)b_0^n$, zero on words containing $b$. The vector $(1,b_0,b_0^2)$ is a top eigenvector, $b_0^6=1/4$, and $\Snorm{H_{h_*}}=2$. Its normalized infinite kernel has squared overlap $3/4$ with the normalized top eigenvector of $K$. The difference of the two rank-one terms of size two has norm $2\sqrt{1-3/4}=1$; the orthogonal eigenvectors of eigenvalues $1,-1$ are unchanged. Thus $\Snorm{H_h-H_{h_*}}=1$, which is optimal by the singular-value bound. Both targets have generated algebra $\mathbb F[J]\simeq\mathbb F[t]/(t^3)$, but the letter $b$ is marked by $J^2$ in one and zero in the other.
\end{proof}
